\chapter{Introduction: Most Startups Fail; How to Improve Your Odds}

\lectureoverview{Joseph Hadzima opens the course by refusing to begin with doctrine. He begins with motive, then with uncertainty, then with odds, and only after that with equations and frameworks. These notes follow his MIT OpenCourseWare lecture in that order, preserving the way the talk raises a local obstacle before it resolves it. The chapter is curated by LazyingArt LLC from the source lecture, and the aim is practical rather than mythic: how we improve our odds without pretending that the odds are good.}

\section{Why We Are Here}

We begin where the lecture begins: with a room full of mixed intentions. Some people are present because they have never taken a formal entrepreneurship class and want only to dip a toe in the water. Some are present because they have seen a product or process that ``sucks'' and want to make it better. Some want to bring an invention to life. Some want to change the world. Some are attracted by the glamour, status, and money that entrepreneurship seems to promise.

Hadzima does not moralize about these motives. He sorts them, and then he turns them into work. The Dropbox anecdote is typical. A student, unable to reach his files while traveling, does not merely complain about the inconvenience; he converts the irritation into a venture thesis. That is the lecture's first serious move. The ordinary annoyance is not yet a business, but it is already a candidate for disciplined investigation.

From there the lecture broadens into a long inventory of first questions. The effect is important. We are not yet given a framework for judgment; first we are made to feel how many moving pieces there are.

\begin{enumerate}
  \item What problem are we solving?
  \item Does the proposed solution actually answer that problem?
  \item Who cares, specifically, and who is the customer?
  \item How will the venture make money?
  \item How long will it take to reach the market, what will it cost, and what resources will be required?
  \item What legal entity, if any, should be formed?
  \item How can the idea be protected, if protection is even the right objective?
  \item Will cofounders be needed, and how will the relationship be structured?
  \item How will the entrepreneur negotiate with employees, consultants, advisors, customers, and investors?
  \item How do we discover what we do not yet know?
\end{enumerate}

The lecture's early rhythm is therefore interrogative. We do not begin with answers. We begin by learning to notice the right questions.

\section{Why Do This If Most Startups Fail?}

Once the room has admitted its motives, the lecture inserts resistance. Why do any of this if the work is painful and the odds are poor? Hadzima sharpens the question with a contrast between startup glamour and startup reality. The glamorous image is fame and fortune. The reality, illustrated with the Nvidia anecdote, is decades of strain, embarrassment, vulnerability, and a continuous encounter with things that may go wrong.

The lecture then gives the base rates in deliberately blunt form:
\begin{align}
\Pr(\text{startup success}) &\approx 0.10 \approx 10\%, \\
\Pr(\text{startup failure}) &\approx 0.90 \approx 90\%, \\
\Pr(\text{top-tier VC success per investment}) &\approx \frac{3}{10} = 0.300.
\end{align}

The baseball analogy matters because it converts an abstract venture-capital ratio into a familiar benchmark. A venture capitalist who succeeds three times out of ten is, by the lecture's standard, a hero in roughly the way a baseball player batting .300 is a hero. The point is not that entrepreneurship is like baseball in every respect. The point is that the failure rate is not an embarrassment added from the outside; it is part of the subject.

And yet the lecture does not stop there. It also gives a local MIT counterweight, not as a proof that entrepreneurship is easy here, but as evidence that success, when it happens, can be very large:
\begin{equation}
\frac{26{,}000}{120{,}000} \approx 0.217 \approx \frac{1}{5}.
\end{equation}
That is the Kauffman-style ratio quoted in the lecture: roughly one active MIT-founded company for every five living alumni at the time of the report. The same framing attributes annual revenues on the order of \(\$2\) trillion to those companies and says that, aggregated, they would amount to roughly the eleventh largest economy in the world.

So the course does not deny the odds. It accepts them and then states its narrower objective: not to abolish risk, but to improve preparedness and thereby improve the probability of landing in the successful minority.

\section{Why This Course, and Who Is In The Room?}

After the statistics, Hadzima rebuilds confidence by explaining what kind of course this is. It is not a theory course about entrepreneurship. It is a doing course, taught by people who either are doing the work or have done it. Its origin story reinforces the same point. Students effectively created demand for a January course on how to start a company; the initial design was assembled quickly from practitioners willing to explain what they wished someone had told them when they started. That story is not ornamental. It establishes the course itself as a small entrepreneurial response to expressed demand.

Then the frame narrows from the course to the room. The audience is heterogeneous: students from engineering, science, architecture and planning, management and economics, and other areas; non-student participants such as alumni, staff, and outside practitioners; and a wide spread of venture interests. The slide matters here because the equation is not floating in isolation. It is staged against the problem of teaching one course to people with very different backgrounds and expectations.

\begin{figure}[H]
\centering

\includegraphics[width=0.92\linewidth]{lecture_01_figure_02.png}
\caption{Audience mix and the course happiness equation.}
\end{figure}

The slide writes the quotient in a stylized stacked form. We standardize it in ordinary mathematical notation:
\begin{equation}
H=\frac{R}{E}.
\end{equation}

\begin{definition}[The Happiness Equation]
In the lecture's notation,
\[
H=\frac{R}{E},
\]
where \(H\) is happiness, \(R\) is reality, and \(E\) is expectations.
\end{definition}

\subsection{Question \& Answer}

\paragraph{Question.} What are \(R\) and \(E\) in \(H=\frac{R}{E}\)?

\paragraph{Answer.} \(R\) is reality and \(E\) is expectations. The lecture briefly records a revealing false start: some Sloan students first propose \(R=\) revenue and \(E=\) expenses. Hadzima accepts that as a subset, then generalizes the equation. He is not giving us an accounting identity. He is giving us a broader rule for interpreting satisfaction, disappointment, and overpromising.

From the equation we can extract the lecture's simple logic:
\begin{align}
H &= \frac{R}{E}, \\
H>1 &\iff R>E, \\
H=1 &\iff R=E, \\
H<1 &\iff R<E.
\end{align}

This is the mathematical content behind the lecture's verbal contrast. If we underpromise and overdeliver, then \(R>E\), and the ratio exceeds one. If we overpromise and underdeliver, then \(R<E\), and the ratio falls below one.

The lecture also says that the goal is to maximize \(H\). We should read that as pedagogical rhetoric, not as a formal constrained optimization problem. Still, it is useful rhetoric. The course wants to raise reality rather than merely manipulate expectations, and the later warning about unsupported hype will return to exactly this point.

\begin{remark}
The lecture jokes that if expectations were driven to zero one could become infinitely happy. We should not literalize that joke into a usable model. The serious content is comparative: inflated expectations without corresponding reality reduce the ratio.
\end{remark}

\section{Frameworks for Judging a Venture}

Once the room has been calibrated, the lecture begins compressing entrepreneurship into usable filters. The first compression is the simplest.

A venture must create value. If it does not create value, it does not have anything worth discussing. But creation alone is not enough. It must also capture enough of that value to continue. In a for-profit venture, that eventually means profit and cash flow. In a nonprofit venture, it may mean enough continuing support and resource inflow to keep operating. The distinction matters because Hadzima refuses to equate entrepreneurship with a single legal or financial form. Business, nonprofit, and even governmental ventures all fall under the course's umbrella, but all must answer the same structural question: how is created value preserved well enough that the work can continue?

The next compression is the ``three whys,'' followed by the risk question that appears only after the first three have begun to work.

\begin{enumerate}
  \item \textbf{Why this?} Why is this problem worth solving, this market worth entering, or this idea worth backing?
  \item \textbf{Why now?} Why is this the right moment rather than five years earlier or later?
  \item \textbf{Why this team?} Why should we believe these people will win?
  \item \textbf{Why won't this work?} What are the failure modes, and how will they be de-risked?
\end{enumerate}

The first why is illustrated by Amy Smith's water-incubator pitch, which is a model of compression. The lecture's force here comes from the numerical opening: \(1.9\) billion people lacking access to clean water. That number is not decoration. It establishes scale before the mechanism arrives. The mechanism is then narrow and concrete: contamination testing requires incubation, existing incubators require electricity, the affected population often lacks electricity, and therefore a non-electric incubator attacks a real bottleneck. The reasoning is sequential and the pitch works because each step narrows the gap opened by the first statistic.

The second why, why now, is framed as convergence. The lecture's FAA example shows that a good technical solution can become fundable when institutional demand appears at the right moment. The same idea will later return in the timing stories about 3D printing, online shopping, and fusion.

The third why, why this team, is treated with healthy suspicion. Prior experience can help, but it is not magic. Even a team assembled from the technical founders of previously successful computer companies may fail. Experience can attract capital; it does not guarantee alignment or victory.

The fourth question, ``Why won't this work?'', is where the lecture begins its de-risking logic. Hadzima's point is that one usually does not get to this question until the first three have generated real interest. Once there, however, the question is essential. We are no longer asking whether the story sounds good. We are testing whether it can survive contact with reality.

\section{Ideas, Execution, Timing, and People}

The lecture's next major compression is the four-part success model: ideas, execution, timing, and people. Here again the order matters. First the framework is displayed in general form. Only afterward does the lecture populate the lower graph and turn the picture into a dynamic argument about team balance.

\begin{figure}[H]
\centering

\includegraphics[width=0.88\linewidth]{lecture_01_figure_03.png}
\caption{Four startup factors above an empty importance-time graph.}
\end{figure}

\begin{figure}[H]
\centering
\begin{tikzpicture}[x=1cm,y=1cm,>=Latex,font=\small]
  \def\r{1.55}
  \coordinate (c1) at (-1.25,7.05);
  \coordinate (c2) at (1.25,7.05);
  \coordinate (c3) at (-1.25,4.65);
  \coordinate (c4) at (1.25,4.65);

  \draw[thick,blue!70!black] (c1) circle (\r);
  \draw[thick,blue!70!black] (c2) circle (\r);
  \draw[thick,blue!70!black] (c3) circle (\r);
  \draw[thick,blue!70!black] (c4) circle (\r);

  \node[anchor=east,font=\bfseries\small,text=blue!70!black] at (-3.05,7.05) {IDEAS};
  \node[anchor=west,font=\bfseries\small,text=blue!70!black] at (3.05,7.05) {EXECUTION};
  \node[anchor=east,font=\bfseries\small,text=blue!70!black] at (-3.05,4.65) {TIMING};
  \node[anchor=west,font=\bfseries\small,text=blue!70!black] at (3.05,4.65) {PEOPLE};

  \draw[fill=orange!18] (-3.55,0) rectangle (3.55,3.15);
  \draw[thick] (-2.55,0.55) -- (-2.55,2.65);
  \draw[thick] (-2.55,0.55) -- (2.55,0.55);
  \node[align=center,anchor=west] at (-3.38,2.15) {RELATIVE\\IMPORTANCE};
  \node[anchor=west] at (2.03,0.10) {TIME};
\end{tikzpicture}
\caption{Pocket redraw of the four-factor framework before the lower graph is interpreted.}
\end{figure}

\begin{remark}
Although the lecture verbally introduces this lower panel as a ``curve graph,'' the visible slide at this stage contains only a scaffold: axes, labels, and empty space. That emptiness is part of the narrative. The dynamics have not yet been stated.
\end{remark}

\subsection{Ideas}

Ideas are abundant, especially around MIT. That is not enough. The lecture immediately asks the harder questions. Valuable to whom? How much value? Can the idea be made durable, or is it easily copied? This is where the earlier create-value logic returns in a sharper form. The entrepreneur does not merely need an idea; the entrepreneur needs an idea whose value can survive contact with a market and with competitors.

\subsection{Execution}

Execution is where attractive ideas begin to die. Hadzima quotes the familiar Edison line, ``Vision without execution is hallucination,'' and then makes the point concrete with investments he declined.

Zipcar looked conceptually important: a serious change in how transportation might be consumed. But the execution burden, especially around parking and city-by-city operational structure, was for him unclear enough to justify passing. The precursor-to-eBay case has the same form. The concept was visible early, but the trust structure and fraud problem were not yet convincing. In both cases the lecture is careful: passing on a venture is not the same as proving it impossible. It is sometimes an admission that the execution problem is not yet intelligible.

This is also why \textit{mens et manus} appears at this point in the lecture. Ideas and execution belong together, but even their conjunction is not sufficient.

\subsection{Timing}

Timing is treated with unusual severity. Hadzima says plainly that he has lost more money and time by being ahead of the curve than by many other mistakes. The examples matter because each isolates a different timing mechanism.

3D printing took decades because the surrounding technologies had to mature. Prodigy spent heavily to build online shopping before the browser and user environment were ready. Fusion is riskier still because it does not work incrementally: all of it must work together. A venture whose value appears only when every component is simultaneously complete carries a timing burden of a different order.

So ``Why now?'' is not a ceremonial question. It is a test of whether the world has arrived at the same moment as the idea.

\subsection{People}

People, in Hadzima's experience, are the single biggest source of venture failure. He makes this point first with the Cheshire Cat story: if we do not know where we want to go, then route choice is meaningless. He then sharpens it with the three-founder example in which one founder wants an industry standard, another wants a fast public company, and a third wants something more collegial and less aggressive. The ideas are not individually absurd. The problem is that they do not define the same journey.

A second people problem appears when teams are funded before they understand what they are doing. Then the pivot is internal rather than customer-facing. The lecture's football-to-basketball analogy is useful here. If the team, stadium, and compensation structure have all been assembled for one game, and the market reveals that it wants another, the crisis is no longer purely strategic. It is organizational.

The E-Ink example adds one more point: not knowing what we do not know is itself a people risk. A team that is certain too early, especially about capital needs, is difficult to work with because it has confused confidence with understanding.

\subsection{Question \& Answer}

\paragraph{Question.} How do technical and business contributions change over time inside a founding team?

\paragraph{Answer.} The lecture turns this into a qualitative time-dependent graph. At the beginning, the technical side dominates because until something works technically, there is no venture to sell, finance, or scale. As time passes, the business side grows in relative importance because market development, staffing, financing, negotiation, and organizational structure become progressively more central. The graph is therefore not a measurement device. It is a way of making a team argument visible.

\begin{equation}
I_{\text{technical}}(t)\downarrow,\qquad I_{\text{business}}(t)\uparrow.
\end{equation}

\begin{figure}[H]
\centering

\includegraphics[width=0.88\linewidth]{lecture_01_figure_04.png}
\caption{Technical importance falls as business importance rises.}
\end{figure}

\begin{figure}[H]
\centering
\begin{tikzpicture}[x=1cm,y=1cm,>=Latex,font=\small]
  \def\r{1.55}
  \coordinate (c1) at (-1.25,7.05);
  \coordinate (c2) at (1.25,7.05);
  \coordinate (c3) at (-1.25,4.65);
  \coordinate (c4) at (1.25,4.65);

  \draw[thick,blue!70!black] (c1) circle (\r);
  \draw[thick,blue!70!black] (c2) circle (\r);
  \draw[thick,blue!70!black] (c3) circle (\r);
  \draw[thick,blue!70!black] (c4) circle (\r);

  \node[anchor=east,font=\bfseries\small,text=blue!70!black] at (-3.05,7.05) {IDEAS};
  \node[anchor=west,font=\bfseries\small,text=blue!70!black] at (3.05,7.05) {EXECUTION};
  \node[anchor=east,font=\bfseries\small,text=blue!70!black] at (-3.05,4.65) {TIMING};
  \node[anchor=west,font=\bfseries\small,text=blue!70!black] at (3.05,4.65) {PEOPLE};

  \draw[fill=orange!18] (-3.55,0) rectangle (3.55,3.15);
  \draw[thick] (-2.55,0.55) -- (-2.55,2.65);
  \draw[thick] (-2.55,0.55) -- (2.55,0.55);
  \node[align=center,anchor=west] at (-3.38,2.15) {RELATIVE\\IMPORTANCE};
  \node[anchor=west] at (2.03,0.10) {TIME};

  \draw[->,very thick,red!85!black] (-1.10,2.55) -- (1.55,1.55);
  \node[anchor=west,font=\bfseries\small,text=red!85!black] at (-2.05,2.72) {TECHNICAL};

  \draw[->,very thick,blue!85!black] (-1.30,1.02) -- (1.48,1.78);
  \node[anchor=west,font=\bfseries\small,text=blue!85!black] at (-2.05,0.86) {BUSINESS};
\end{tikzpicture}
\caption{Pocket redraw of the lecture's qualitative team-balance graph.}
\end{figure}

\begin{remark}
The graph is schematic. The arrows are straight, unscaled, and not tied to a marked crossing time. The lecture explicitly says that whether the lines cross is not the point. The point is that role importance changes, and teams that ignore that fact will fight about structure, equity, and legitimacy.
\end{remark}

\subsection{Four Cases}

The lecture then runs the framework against four ventures.

\begin{enumerate}
  \item \textbf{SpeechCo.} The lecture presents this as a case where the four factors align: the idea was strong, the timing was right because compute power had arrived, execution improved by restricting the system to a domain of language use, and the people had spent enough time together to form a working team.
  \item \textbf{VideoCo.} Here Hadzima initially doubted execution because the founders appeared unable to settle even on a company name. The payoff is memorable: they were not merely naming the company but choosing a strategy, and the eventual DIVA/Avid connection revealed that the supposed indecision had a directional logic.
  \item \textbf{HIV Co.} This case had a strong idea, strong timing, and substantial scientific credibility. The weakness was people. Cofounder conflict drained the venture even after powerful backers and advisers had entered the scene.
  \item \textbf{NanoCo.} This is the clean timing failure. The product, NanoFuel, had a serious technical claim and an environmental advantage, but market timing broke the venture's path to adoption.
\end{enumerate}

NanoFuel also produces the lecture's clearest explicit back-of-the-envelope calculation:
\begin{align}
p_{\text{diesel}} &\approx \$0.50/\text{gal}, \\
p_{\text{neutral}} &\approx \$2.00/\text{gal}.
\end{align}

At launch,
\begin{equation}
p_{\text{diesel}} < p_{\text{neutral}},
\end{equation}
and therefore the product sits below the lecturer's rough economic-neutral threshold. The argument is intentionally simple, but it is enough for the lecture's point. A technically interesting product may still fail because the world has supplied the wrong price regime at the wrong time.

\section{From Vision to Pitch, Without Losing the Substance}

Only after a long discussion of substance does the lecture turn to communication. That order is essential. We are not being taught to replace venture content with venture theater. We are being taught to compress a real venture into forms that other people can quickly evaluate.

\begin{figure}[H]
\centering
\includegraphics[width=0.82\linewidth]{lecture_01_figure_05.png}
\caption{A supported vision from mission statement to pitch deck.}
\end{figure}

\begin{figure}[H]
\centering
\begin{tikzpicture}[x=1cm,y=1cm,>=Latex,font=\small]
  % pyramid tiers
  \fill[red!65!magenta!35] (0,4.20) -- (0.28,3.55) -- (-0.28,3.55) -- cycle;
  \fill[magenta!55] (-0.28,3.55) -- (-0.54,2.95) -- (0.54,2.95) -- (0.28,3.55) -- cycle;
  \fill[yellow!85!orange!25] (-0.54,2.95) -- (-1.01,1.85) -- (1.01,1.85) -- (0.54,2.95) -- cycle;
  \fill[green!75!black!20] (-1.01,1.85) -- (-1.80,0) -- (1.80,0) -- (1.01,1.85) -- cycle;

  \draw[thick] (0,4.20) -- (-1.80,0) -- (1.80,0) -- cycle;
  \draw[thick] (-0.28,3.55) -- (0.28,3.55);
  \draw[thick] (-0.54,2.95) -- (0.54,2.95);
  \draw[thick] (-1.01,1.85) -- (1.01,1.85);

  % callouts
  \draw[->,thick] (2.55,3.95) -- (0.04,3.95);
  \node[anchor=west] at (2.62,3.95) {Mission/Vision Statement};
  \node[anchor=west,text=red!80!black] at (5.65,3.95) {1 sentence/paragraph};

  \draw[->,thick] (2.55,3.25) -- (0.12,3.25);
  \node[anchor=west] at (2.62,3.25) {Elevator Pitch};

  \draw[->,thick] (2.55,2.35) -- (0.28,2.35);
  \node[anchor=west] at (2.62,2.35) {Executive Summary};

  \draw[->,thick] (2.55,1.05) -- (0.72,1.05);
  \node[anchor=west] at (2.62,1.05) {Pitch Deck};
\end{tikzpicture}
\caption{Pocket redraw of the supported-vision pyramid.}
\end{figure}

The pyramid is hierarchical, not quantitative. The mission statement is the most compressed form. The elevator pitch unfolds it orally. The executive summary adds short written support. The pitch deck adds a fuller argumentative structure. The lecture gives approximate communication scales:
\begin{align}
\text{mission statement} &\sim \text{1 sentence or paragraph}, \\
\text{elevator pitch} &\sim 30\text{ seconds}, \\
\text{executive summary} &\sim 1\text{--}4\text{ pages}, \\
(\text{slides},\text{minutes},\text{font size}) &= (10,20,30).
\end{align}

That last line is Guy Kawasaki's 10/20/30 rule: ten slides, twenty minutes, thirty-point font. The lecture presents it as discipline, not decoration. Compression is useful only if the compressed statement is supported by something real underneath it.

This is where the Theranos warning enters. A venture may get the sizzle right while the support is missing. If that happens, the pyramid is upside down. The lecture's caution can be expressed directly in the language of the happiness equation. If reality is fixed while expectations are inflated, the ratio worsens:
\begin{equation}
R \text{ fixed},\qquad E_2>E_1 \qquad \Longrightarrow \qquad \frac{R}{E_2}<\frac{R}{E_1}.
\end{equation}

So communication is not a separate art laid on top of the venture after the fact. It is the distillation of the venture. If the underlying support is thin, then a brilliant pitch merely increases the size of the disappointment.

\section{Summary}

The lecture's structure is now visible. We begin with motive and uncertainty. We then confront bad odds. We justify the course not by denying those odds, but by arguing that preparation improves them. The room itself is turned into a local equation, \(H=\frac{R}{E}\), and that equation becomes the template for a broader warning against inflated promises. Then the lecture compresses venture judgment into working filters: create value and capture value; why this, why now, why this team, and why won't this work; ideas, execution, timing, and people.

The final lesson is modest and important. Preparation does not remove luck. It makes us more ready when luck matters. That is why the lecture closes without triumphalism. Even deciding that a particular venture is not the right one can count as success, provided we have learned enough not to make the same mistake twice.